% from ega2-1-fr.tex
\oldpage[II]{103}

dans les propositions (\hyperref[II.5.4.2-fr]{5.4.2}) à
(\hyperref[II.5.4.8-fr]{5.4.8}), on peut partout remplacer «~propre~» par
«~universellement fermé~» sans modifier la validité des résultats (et dans
les hypothèses de (\hyperref[II.5.4.3-fr]{5.4.3}),
(\hyperref[II.5.4.5-fr]{5.4.5}), (\hyperref[II.5.4.6-fr]{5.4.6}) et
(\hyperref[II.5.4.8-fr]{5.4.8}), on peut omettre les conditions de
finitude).
\end{remark}

\begin{env}[5.4.10]
\label{II.5.4.10-fr}
Soit $f:X\to Y$ un morphisme de type fini. Nous dirons qu'une partie fermée
$Z$ de $X$ est \emph{propre sur $Y$} (ou \emph{$Y$-propre}, ou \emph{propre
pour $f$}) si la restriction de $f$ à un sous-préschéma fermé de $X$,
d'espace sous-jacent $Z$ (\hyperref[I.5.2.1-fr]{I, 5.2.1}), est propre.
Comme cette restriction est alors séparée, il résulte de
(\hyperref[II.5.4.6-fr]{5.4.6}) et (\hyperref[I.5.5.1-fr]{I, 5.5.1, (vi)})
que la propriété précédente ne dépend pas du sous-préschéma fermé de $X$
ayant $Z$ pour espace sous-jacent. Si $g:X'\to X$ est un morphisme propre,
$g^{-1}(Z)$ est une partie propre de $X'$~: si $T$ est un sous-préschéma de
$X$ ayant pour espace sous-jacent $Z$, il suffit en effet de remarquer que
la restriction de $g$ au sous-préschéma fermé $g^{-1}(T)$ de $X'$ est un
morphisme propre $g^{-1}(T)\to T$ par
(\hyperref[II.5.4.2-fr]{5.4.2, (iii)}), et d'appliquer ensuite
(\hyperref[II.5.4.2-fr]{5.4.2, (ii)}). D'autre part, si $X''$ est un
$Y$-schéma de type fini et $u:X\to X''$ un $Y$-morphisme, $u(Z)$ est une
partie propre de $X''$~; en effet, prenons pour $T$ le sous-préschéma fermé
réduit de $X$ ayant $Z$ pour espace sous-jacent~; la restriction de $f$ à
$T$ étant propre, il en est de même de la restriction de $u$ à $T$
(\hyperref[II.5.4.3-fr]{5.4.3, (i)}), donc $u(Z)$ est fermé dans $X''$~;
soit $T''$ un sous-préschéma fermé de $X''$ ayant $u(Z)$ pour espace
sous-jacent (\hyperref[I.5.2.1-fr]{I, 5.2.1}), de sorte que $u|T$ se
factorise en
\[
  T\xrightarrow{v}T''\xrightarrow{j}X'',
\]
où $j$ est l'injection canonique (\hyperref[I.5.2.2-fr]{I, 5.2.2}) et $v$
est donc propre et surjectif (\hyperref[II.5.4.5-fr]{5.4.5})~; si $g$ est la
restriction à $T''$ du morphisme structural $X''\to Y$, $g$ est séparé et
de type fini, et on a $f|T=g\circ v$~; il résulte donc de
(\hyperref[II.5.4.3-fr]{5.4.3, (ii)}) que $g$ est propre, d'où notre
assertion.
\end{env}

Il résulte en particulier de ces remarques que si $Z$ est une partie
$Y$-propre de $X$, alors~:
\begin{enumerate}
  \item[\rm{1\textsuperscript{o}}] Pour tout sous-préschéma fermé $X'$ de
  $X$, $Z\cap X'$ est une partie $Y$-propre de $X'$.
  \item[\rm{2\textsuperscript{o}}] Si $X$ est un sous-préschéma d'un
  $Y$-schéma de type fini $X''$, $Z$ est aussi une partie $Y$-propre de
  $X''$ (et en particulier est fermée dans $X''$).
\end{enumerate}

\subsection{Morphismes projectifs.}
\label{subsection:II.5.5-fr}

\begin{proposition}[5.5.1]
\label{II.5.5.1-fr}
Soit $X$ un $Y$-préschéma. Les conditions suivantes sont équivalentes~:
\begin{enumerate}
  \item[\rm{a)}] $X$ est $Y$-isomorphe à un sous-préschéma fermé d'un fibré
  projectif $\mathbf{P}(\mathcal{E})$, où $\mathcal{E}$ est un
  $\mathcal{O}_Y$-Module quasi-cohérent de type fini.
  \item[\rm{b)}] Il existe une $\mathcal{O}_Y$-Algèbre graduée
  quasi-cohérente $\mathcal{S}$, telle que $\mathcal{S}_1$ soit de type fini
  et engendre $\mathcal{S}$, et que $X$ soit $Y$-isomorphe à
  $\operatorname{Proj}(\mathcal{S})$.
\end{enumerate}
\end{proposition}

La condition a) entraîne b) en vertu de
(\hyperref[II.3.6.2-fr]{3.6.2, (ii)})~: si $\mathcal{J}$ est un faisceau
gradué quasi-cohérent d'idéaux de $\mathbf{S}(\mathcal{E})$, la
$\mathcal{O}_Y$-Algèbre graduée quasi-cohérente
$\mathcal{S}=\mathbf{S}(\mathcal{E})/\mathcal{J}$ est engendrée par
$\mathcal{S}_1$ et ce dernier, image canonique de $\mathcal{E}$, est un
$\mathcal{O}_Y$-Module de type fini. La condition b) implique a) en vertu de
(\hyperref[II.3.6.2-fr]{3.6.2}) appliqué au cas où
$\mathcal{M}\to\mathcal{S}_1$ est l'application identique.

